\documentclass[10pt,a4paper]{article}
\usepackage{amsmath,amssymb,amsthm,graphicx}
\usepackage{amsmath,amssymb,amsthm,hyperref}
\newcommand{\Av}{\bar A}\newcommand{\Bv}{\bar B}\newcommand{\Cv}{\bar C}
\newcommand{\sv}{\bar\sigma}\newcommand{\tv}{\bar\tau}
\newcommand{\vsub}{\mathrel{\overline{\subseteq}}}
\newcommand{\vin}{\mathrel{\overline{\in}}}
\newcommand{\vcup}{\mathbin{\overline{\cup}}}
\newcommand{\dimn}{\operatorname{dim}}
\newcommand{\dom}{\operatorname{dom}}
\newcommand{\form}{\operatorname{form}}
\newcommand{\pmjudgment}[3]{#1\vdash_k #2:#3}
\newcommand{\pmstrictjudgment}[3]{#1\vdash_S #2:#3}
\newcommand{\orig}[2]{\par\medskip\noindent\textbf{#1 #2.}\ }
\title{Decidability of inhabitation in the bounded parallel intersection type system}
\author{Andrej Dudenhefner, Aleksy Schubert, Jakob Rehof}
\date{doi:10.4230/LIPIcs.FSCD.2026.17; CC BY4.0}
\begin{document}
\maketitle
\noindent\textit{Editorial note.} Complete author upper-bound result transcribed from its official published CC BY4.0 PDF. No author native TeX claim is made. Original PDF and source positions retained.\par\medskip
\noindent\textit{Editorial note.} The general dimension and fixed dimension cases are counted together. Every specialized proof continuation used by the algorithm is retained; examples and unrelated lower bounds/BCD results are uncounted.\par\medskip
\noindent\textit{Editorial note.} Faithful transcription preserves author notation and terse proof steps; no mathematical correction or independent certification performed.\par\medskip
\noindent\textit{Editorial scope.} Complete author material collection for one structural theorem; author retains responsibility for the mathematics.\par\medskip
\noindent\textit{Allowed standard prerequisites (editorial).}\begin{itemize}
\item Strict intersection-type normalization and beta-normal inhabitant reduction. Author cites van Bakel[21] Theorems6.5 and7.5; exact strict-system rules and paper-specific projection Lemma27 retained.
\item Alternating exponential-space equals deterministic double-exponential-time; alternating polynomial-space equals EXPTIME. Standard complexity identities invoked in the complete author Theorems37/38 proof; not independently re-proved.
\end{itemize}\medskip
% SOURCE published pages17:4--17:5. Author notation and definitions retained.
\section*{System R: necessary author definitions}
We denote intersection types by $A,B$, finite sets of intersection types by $\sigma,\tau$, and type atoms by $a,b$.
\orig{Definition}{1 (Intersection Types)}
\[A,B ::= a\mid\sigma\to A\qquad\sigma,\tau::=\{A_1,\ldots,A_n\}\text{ where }n\geq0.\]
Intersection types are a variant of strict types considered by van Bakel [20,21,22].
We denote vectors of types and vectors of sets of types by $\Av$ and $\sv$ respectively. If $\sv=(\sigma_1,\ldots,\sigma_n)$ and $\Av=(A_1,\ldots,A_n)$, then $\sv\Rightarrow\Av=(\sigma_1\to A_1,\ldots,\sigma_n\to A_n)$, where $n\geq0$, and in particular $()\Rightarrow()=()$. We write $\Av_{[i]}$ for $A_i$ and $\sv_{[i]}$ for $\sigma_i$. The dimension of a vector is the number of its coordinates, that is $\dimn(A_1,\ldots,A_n)=n$. We extend relations $\subseteq,\in$, and the operation $\cup$ to vectors.
\orig{Definition}{2 (Vector Extension)}
\begin{align*}
(\sigma_1,\ldots,\sigma_n)\vsub(\tau_1,\ldots,\tau_n)&\quad\text{if }\sigma_1\subseteq\tau_1,\ldots,\sigma_n\subseteq\tau_n,\\
(A_1,\ldots,A_n)\vin(\sigma_1,\ldots,\sigma_n)&\quad\text{if }A_1\in\sigma_1,\ldots,A_n\in\sigma_n,\\
(\sigma_1,\ldots,\sigma_n)\vcup(\tau_1,\ldots,\tau_n)&:=(\sigma_1\cup\tau_1,\ldots,\sigma_n\cup\tau_n),\\
(\{A_1\},\ldots,\{A_n\})\mathord\downarrow&:=(A_1,\ldots,A_n).
\end{align*}
\orig{Definition}{3 (Environment of Dimension $d$)}
An environment $\Sigma$ is a finite set of type assumptions having the shape $x:\sigma$ for distinct term variables. An environment $\Gamma$ of dimension $d$ is a finite set of type assumptions having the shape $x:\sv$ for distinct term variables such that $\dimn(\sv)=d$.
Each environment $\Gamma$ is equipped with its dimension $\dimn(\Gamma)$. In particular, the empty environment of dimension $0$ is distinct from the empty environment of dimension $1$.
We may extend an environment $\Gamma$ by an additional assumption $x:\sv$, written $\Gamma,x:\sv$, where $\dimn(\sv)=\dimn(\Gamma)$ and $x$ does not appear in any assumption in $\Gamma$. The domain of an environment $\Gamma$ is given by $\dom(\Gamma)=\{x\mid(x:\sv)\in\Gamma\text{ and }\sv\neq(\emptyset,\ldots,\emptyset)\}$. When $(x:\sv)\in\Gamma$ we write $\Gamma(x)$ for $\sv$. We identify $\Gamma$ and $\Gamma,x:(\emptyset,\ldots,\emptyset)$, consequently $\Gamma(y)=(\emptyset,\ldots,\emptyset)$ iff $y\notin\dom(\Gamma)$ for any term variable $y$.
Similarly to vectors, we extend the relation $\subseteq$, the operation $\cup$, and coordinate indexing to environments.
\orig{Definition}{4 (Environment Extension)}
Let $\Gamma,\Delta$ be environments such that $\dimn(\Gamma)=\dimn(\Delta)$.
\begin{align*}
\Gamma\subseteq\Delta&\quad\text{if }\Gamma(x)\vsub\Delta(x)\text{ for }x\in\dom(\Gamma),\\
(\Gamma\cup\Delta)(x)&:=\Gamma(x)\vcup\Delta(x),\\
\Gamma_{[i]}&:=\{x:\sv_{[i]}\mid(x:\sv)\in\Gamma\}.
\end{align*}
The original feature of the system presented in this paper is the management of flow of types between vectors by means of relations. We extend notation of previous work [10] from surjective functions to binary relations, and define type vector transformations as follows.
\orig{Definition}{5 (Type Vector Transformation)}
Let $R\subseteq\{1,\ldots,n\}\times\{1,\ldots,m\}$ be a binary relation where $n,m\geq0$.
\begin{align*}
R((A_1,\ldots,A_n))&:=(\tau_1,\ldots,\tau_m)\text{ where }\tau_i=\{A_j\mid j\,R\,i\},\\
R((\sigma_1,\ldots,\sigma_n))&:=(\tau_1,\ldots,\tau_m)\text{ where }\tau_i=\bigcup\{\sigma_j\mid j\,R\,i\},\\
R(\{x_1:\sv^1,\ldots,x_l:\sv^l\})&:=\{x_1:R(\sv^1),\ldots,x_l:R(\sv^l)\}.
\end{align*}
\orig{Definition}{7}
Let $R\subseteq\{1,\ldots,n\}\times\{1,\ldots,m\}$ be a relation.
\begin{itemize}
\item $R$ is left-total when for each $i\in\{1,\ldots,n\}$ there is a $j$ such that $i\,R\,j$.
\item $R$ is right-total when for each $j\in\{1,\ldots,m\}$ there is an $i$ such that $i\,R\,j$.
\item $R$ is left-unique when for each $i\in\{1,\ldots,n\}$ there is at most one $j$ such that $i\,R\,j$.
\item $R$ is right-unique when for each $j\in\{1,\ldots,m\}$ there is at most one $i$ such that $i\,R\,j$.
\end{itemize}
If a relation $R$ is right-unique and right-total, then the vector of types $R(\Bv)\mathord\downarrow$ is well-defined. Primary properties of vector transformations are listed in the following Lemma 8.
\orig{Lemma}{8 (Properties of Vector Transformations)}
\begin{enumerate}
\item For two binary relations $R,R'\subseteq\{1,\ldots,n\}\times\{1,\ldots,m\}$ it holds that if $R\subseteq R'$, then for all $\sv$ such that $\dimn(\sv)=n$ we have $R(\sv)\vsub R'(\sv)$.
\item If $R(\sv)\vsub R'(\Av)$ and $R$ is left-total, then there exists $R''$ such that $R''(\Av)=\sv$.
\item If $\Av$ has distinct coordinates, $R''(\Av)=\sv$, $R(\sv)\vsub R'(\Av)$, and $R$ is left-total, then $R\circ R''\subseteq R'$.
\item $R(\Gamma_1\cup\Gamma_2)=R(\Gamma_1)\cup R(\Gamma_2)$.
\item If $R$ is right-unique and right-total, and $\Av\vin\sv$ then $R(\Av)\mathord\downarrow\vin R(\sv)$.
\item If $R$ is right-unique and right-total, then $R(\sv\Rightarrow\Av)\mathord\downarrow=R(\sv)\Rightarrow R(\Av)\mathord\downarrow$.
\item If $R(\sv)\vsub\tv$ and $R$ is left-total, then there exists $\sv'$ such that $\sv\vsub\sv'$, $R(\sv')\vsub\tv$, and for all $i\in\{1,\ldots,\dimn(\sv')\}$ there exists $j\in\{1,\ldots,\dimn(\tv)\}$ such that $\sv'_{[i]}=\tv_{[j]}$.
\end{enumerate}
\begin{proof}
\begin{enumerate}
\item The claim follows from $R(\sv)_{[i]}=\bigcup\{\sv_{[j]}\mid j\,R\,i\}\subseteq\bigcup\{\sv_{[j]}\mid j\,R'\,i\}=R'(\sv)_{[i]}$.
\item Consider $R''=\{(i,j)\mid\exists k.i\,R'\,k\text{ and }j\,R\,k\text{ and }\Av_{[i]}\in\sv_{[j]}\}$. We have $R''(\Av)_{[j]}=\{\Av_{[i]}\mid i\,R''\,j\}=\{\Av_{[i]}\mid\exists k.i\,R'\,k\text{ and }j\,R\,k\text{ and }\Av_{[i]}\in\sv_{[j]}\}$. This set is included in $\sv_{[j]}$ as for each its element $\Av_{[i]}$ we have $\Av_{[i]}\in\sv_{[j]}$. Similarly, $\sv_{[j]}\subseteq R(\sv)_{[k]}$ for some $k$ as $R$ is left-total. As $R(\sv)_{[k]}\subseteq R'(\Av)_{[k]}$, each element of $\sv_{[j]}$ must be equal to some $\Av_{[i]}$ such that $i\,R'\,k$.
\item We have $(R\circ R'')(\Av)=R(\sv)\vsub R'(\Av)$. Since $\Av$ has distinct coordinates, we obtain the claim.
\item Immediate from the definition.
\item Since $R$ is right-unique and right-total, for each $i\in\{1,\ldots,\dimn(R(\Av)\mathord\downarrow)\}$ there is exactly one $j$ such that $j\,R\,i$, and we have $(R(\Av)\mathord\downarrow)_{[i]}=\Av_{[j]}\in\sv_{[j]}\subseteq R(\sv)_{[i]}$.
\item Since $R$ is right-unique and right-total, we have $(R(\sv\Rightarrow\Av)\mathord\downarrow)_{[i]}=\sv_{[j]}\to\Av_{[j]}$ where $j\,R\,i$ for exactly one $j$. Similarly, $R(\sv)_{[i]}=\sv_{[j]}$ and $(R(\Av)\mathord\downarrow)_{[i]}=\Av_{[j]}$ where $j\,R\,i$. Consequently, we obtain $R(\sv\Rightarrow\Av)\mathord\downarrow=R(\sv)\Rightarrow R(\Av)\mathord\downarrow$.
\item Since $R$ is left-total, for all $i\in\{1,\ldots,\dimn(\sv)\}$ there exists $j\in\{1,\ldots,\dimn(\tv)\}$ such that $i\,R\,j$, for which we set $\sv'_{[i]}:=\tv_{[j]}$, satisfying the required properties.
\end{enumerate}
\end{proof}
Judgments in the following system R with dimension bound $k$ are of shape $\pmjudgment\Gamma t\Av$ such that $\dimn(\Gamma)=\dimn(\Av)$.
\orig{Definition}{9 (System R with Dimension Bound $k$)}
\[\frac{\Av\vin\sv\qquad\dimn(\Gamma)=\dimn(\Av)\leq k}{\pmjudgment{\Gamma,x:\sv}{x}{\Av}}\ (\mathrm{Ax})\qquad
\frac{\pmjudgment{\Gamma,x:\sv}{t}{\Av}}{\pmjudgment\Gamma{\lambda x.t}{\sv\Rightarrow\Av}}\ (\Rightarrow\mathrm I)\]
\[\frac{\pmjudgment\Gamma t{R(\Av)\Rightarrow\Bv}\qquad\pmjudgment\Delta u\Av\qquad R\subseteq\{1,\ldots,\dimn(\Delta)\}\times\{1,\ldots,\dimn(\Gamma)\}}{\pmjudgment{\Gamma\cup R(\Delta)}{t\,u}{\Bv}}\ (\Rightarrow\mathrm E)\]
We may leave out the actual bound, if it is immaterial.
\orig{Definition}{10 (System R)} $\Gamma\vdash t:\Av$ means $\pmjudgment\Gamma t\Av$ for some $k$.
% SOURCE published page17:6--17:7. Selected support; examples are not counted.
\orig{Lemma}{11} The following rule is derivable: $\displaystyle\frac{}{\pmjudgment\emptyset t{()}}\ (\mathrm{nil})$.
\begin{proof}
We proceed by induction on $t$.\par
\emph{Case $t=x$:} Due to $\emptyset,x:()=\emptyset$, by rule (Ax) we have $\pmjudgment\emptyset x{()}$.\par
\emph{Case $t=\lambda x.u$:} By the induction hypothesis we have $\pmjudgment\emptyset u{()}$. Due to $\emptyset,x:()=\emptyset$ and $()\Rightarrow()=()$, by rule $(\Rightarrow\mathrm I)$ we have $\pmjudgment\emptyset{\lambda x.u}{()}$.\par
\emph{Case $t=u\,v$:} By the induction hypothesis we have $\pmjudgment\emptyset u{()}$ and $\pmjudgment\emptyset v{()}$. For $R=\emptyset$ we have $()=R(())\Rightarrow()$ and $\emptyset\cup R(\emptyset)=\emptyset$. Therefore, by rule $(\Rightarrow\mathrm E)$ we have $\pmjudgment\emptyset{u\,v}{()}$.
\end{proof}
\orig{Lemma}{14 (Dimensional Restriction)} Any vector occurring in a derivation of $\pmjudgment\Gamma t\Av$ has at most $k$ coordinates.
\begin{proof}Induction on $t$.\end{proof}
\orig{Lemma}{15 (Dimension Weakening)} If $\pmjudgment\Gamma t\Av$ and $k\leq m$ then $\Gamma\vdash_m t:\Av$.
\begin{proof}Induction on $t$. For rule (Ax) the dimension restriction $n\leq k$ is replaced by $n\leq m$, which holds given $k\leq m$.\end{proof}
\orig{Lemma}{16 (Weakening)} If $\pmjudgment\Gamma t\Av$ and $\Gamma\subseteq\Delta$, then $\pmjudgment\Delta t\Av$.
\begin{proof}
Assuming $\pmjudgment\Gamma t\Av$ and $\Gamma\subseteq\Delta$, we proceed by induction on $t$.\par
\emph{Case $t=x$:} Due to $\Av\vin\Gamma(x)\vsub\Delta(x)$ we obtain the claim by rule (Ax).\par
\emph{Case $t=\lambda x.u$:} For some $\sv$ and $\Bv$ we have $\Av=\sv\Rightarrow\Bv$, and by the induction hypothesis we have $\pmjudgment{\Delta,x:\sv}{u}{\Bv}$. By rule $(\Rightarrow\mathrm I)$ we obtain the claim.\par
\emph{Case $t=u\,v$:} For some $\Delta',\Gamma',\Bv$, and $R$ such that $\Gamma=\Gamma'\cup R(\Delta')$ we have $\pmjudgment{\Gamma'}u{R(\Bv)\Rightarrow\Av}$ and $\pmjudgment{\Delta'}v\Bv$. Due to $\Gamma'\subseteq\Gamma\subseteq\Delta$ by the induction hypothesis we have $\pmjudgment\Delta u{R(\Bv)\Rightarrow\Av}$. By rule $(\Rightarrow\mathrm E)$ we have $\pmjudgment{\Delta\cup R(\Delta')}{u\,v}{\Av}$. Due to $R(\Delta')\subseteq\Gamma\subseteq\Delta$ we have $\Delta\cup R(\Delta')=\Delta$, obtaining the claim.
\end{proof}
A key technical insight for system R is the following Lemma 17, which allows us to permute, remove, or copy vector coordinates in derivations.
\orig{Lemma}{17} If $\pmjudgment\Gamma t\Av$ and $R\subseteq\{1,\ldots,\dimn(\Av)\}\times\{1,\ldots,m\}$ is right-unique and right-total, then $R(\Gamma)\vdash_{\max(k,m)}t:R(\Av)\mathord\downarrow$.
\begin{proof}
Assuming $\pmjudgment\Gamma t\Av$ and that $R$ is right-unique and right-total, we proceed by induction on $t$. Let $n=\max(k,m)$.\par
\emph{Case $t=x$:} We have $\Av\vin\Gamma(x)$. By Lemma 8(5), we have $R(\Av)\mathord\downarrow\vin R(\Gamma(x))$, which by rule (Ax) implies $R(\Gamma)\vdash_n x:R(\Av)\mathord\downarrow$.\par
\emph{Case $t=\lambda x.u$:} For some $\sv$ and $\Bv$ we have $\Av=\sv\Rightarrow\Bv$, and by the induction hypothesis we have $R(\Gamma),x:R(\sv)\vdash_n u:R(\Bv)\mathord\downarrow$. By Lemma 8(6) we have $R(\sv)\Rightarrow R(\Bv)\mathord\downarrow=R(\sv\Rightarrow\Bv)\mathord\downarrow=R(\Av)\mathord\downarrow$. By rule $(\Rightarrow\mathrm I)$ we obtain $R(\Gamma)\vdash_n\lambda x.u:R(\Av)\mathord\downarrow$.\par
\emph{Case $t=u\,v$:} For some $\Delta,R'$, and $\Bv$ we have $\pmjudgment\Delta v\Bv$, and weakening the environment by Lemma 16, we have $\Gamma\vdash_n u:R'(\Bv)\Rightarrow\Av$. By the induction hypothesis we obtain $R(\Gamma)\vdash_n u:R(R'(\Bv)\Rightarrow\Av)\mathord\downarrow$. By Lemma 15 we have $\Delta\vdash_n v:\Bv$. Let $R''=R\circ R'$. By Lemma 8(6) we have $R(R'(\Bv)\Rightarrow\Av)\mathord\downarrow=R''(\Bv)\Rightarrow R(\Av)\mathord\downarrow$. By Lemma 8(4) we have $R(\Gamma)\cup R''(\Delta)=R(\Gamma\cup R'(\Delta))$. By rule $(\Rightarrow\mathrm E)$ we obtain $R(\Gamma\cup R'(\Delta))\vdash_n u\,v:R(\Av)\mathord\downarrow$.
\end{proof}
% SOURCE published pages17:8--17:9. Complete author substitution and subject reduction.
\orig{Lemma}{23 (Substitution Lemma)} If $\pmjudgment{\Gamma,x:R(\Av)}t\Bv$ and $\pmjudgment\Delta u\Av$, then $\pmjudgment{\Gamma\cup R(\Delta)}{t[x:=u]}\Bv$.
\begin{proof}
Assuming $\pmjudgment{\Gamma,x:R(\Av)}t\Bv$ and $\pmjudgment\Delta u\Av$, we proceed by induction on $t$. We also note here that $\pmjudgment{\Gamma,x:R(\Av)}t\Bv$ implies that $R\subseteq\{1,\ldots,\dimn(\Av)\}\times\{1,\ldots,m\}$ where $m\leq k$. W.l.o.g. $\Av$ has distinct coordinates as we can remove duplicate ones by Lemma 17.\par
\emph{Case $t=x$:} We have $\Bv_{[i]}\in R(\Av)_{[i]}$ for $i\in\{1,\ldots,\dimn(\Bv)\}$. Therefore, there is a right-unique $R'\subseteq R$ such that $R'(\Av)=\Bv$. By Lemma 17 we have $\pmjudgment{R'(\Delta)}u{R'(\Av)\mathord\downarrow}$. Due to $R'\subseteq R$, we obtain $\pmjudgment{\Gamma\cup R(\Delta)}u\Bv$ by properties of vector transformations and the weakening lemma (cf. Lemma 8(1) and Lemma 16).\par
\emph{Case $t=y\neq x$:} We have $\pmjudgment\Gamma y\Bv$ and the claim follows by the weakening lemma (cf. Lemma 16).\par
\emph{Case $t=\lambda y.v$:} We may assume due to $\alpha$-conversion that $y\neq x$ and $y\notin\operatorname{var}(u)$. For some $\sv$ and $\Cv$ we have $\Bv=\sv\Rightarrow\Cv$ and $\pmjudgment{\Gamma,x:R(\Av),y:\sv}{v}{\Cv}$. By the induction hypothesis we obtain $\pmjudgment{\Gamma\cup R(\Delta),y:\sv}{v[x:=u]}\Cv$. The claim follows by rule $(\Rightarrow\mathrm I)$.\par
\emph{Case $t=v\,w$:} For some $\Gamma',\Delta',R',\sv,\tv$, and $\Cv$ such that $\Gamma=\Gamma'\cup R'(\Delta')$ and $R(\Av)=\tv\vcup R'(\sv)$ we have $\pmjudgment{\Gamma',x:\tv}v{R'(\Cv)\Rightarrow\Bv}$ and $\pmjudgment{\Delta',x:\sv}w\Cv$.
Note that $R'\subseteq\{1,\ldots,\dimn(\Cv)\}\times\{1,\ldots,\dimn(\Bv)\}$. W.l.o.g. $R'$ is left-total. In order to obtain the claim by rule $(\Rightarrow\mathrm E)$ we consider the corresponding premises.
\begin{itemize}
\item Due to $\tv\vsub R(\Av)$, by the weakening lemma (cf. Lemma 16) and the induction hypothesis we have $\pmjudgment{\Gamma'\cup R(\Delta)}{v[x:=u]}{R'(\Cv)\Rightarrow\Bv}$.
\item Due to $R'(\sv)\vsub R(\Av)$ and since $R'$ is left-total and $\Av$ has distinct coordinates, by Lemma 8(2) and Lemma 8(3) there exists a $R''$ such that $R'\circ R''\subseteq R$ and $R''(\Av)=\sv$. By the induction hypothesis we have $\pmjudgment{\Delta'\cup R''(\Delta)}{w[x:=u]}\Cv$.
\end{itemize}
By rule $(\Rightarrow\mathrm E)$ we have $\pmjudgment{\Gamma'\cup R(\Delta)\cup R'(\Delta'\cup R''(\Delta))}{(v\,w)[x:=u]}\Bv$. By the weakening lemma (cf. Lemma 16), it suffices to show that $\Gamma'\cup R(\Delta)\cup R'(\Delta')\cup R'(R''(\Delta))\subseteq\Gamma\cup R(\Delta)$. This follows from $\Gamma=\Gamma'\cup R'(\Delta')$ and $R'\circ R''\subseteq R$.
\end{proof}
\orig{Lemma}{24 (Subject Reduction)} If $\pmjudgment\Gamma t\Bv$ and $t\to_\beta u$ then $\pmjudgment\Gamma u\Bv$.
\begin{proof}
The proof is by induction over $t$.\par
\emph{Case $t=x$:} This case is impossible since $x$ is in $\beta$-normal form.\par
\emph{Case $t=\lambda x.v$:} In this case, reduction actually has the form $t=\lambda x.v\to_\beta\lambda x.v'=u$ where $v\to_\beta v'$. The only way to derive type for $\lambda x.v$ is to use rule $(\Rightarrow\mathrm I)$. This means that $\Bv=\sv\Rightarrow\Av$ for some $\sv$ and $\Av$. From that we get a derivation for $\pmjudgment{\Gamma,x:\sv}v\Av$. We can now use the induction hypothesis to get $\pmjudgment{\Gamma,x:\sv}{v'}\Av$ and then rule $(\Rightarrow\mathrm I)$ to get derivability of $\pmjudgment\Gamma{v'}{\sv\Rightarrow\Av}$, which is the required conclusion.\par
\emph{Case $t=v\,w$:} In this case reduction has two forms. In the first of them, we reduce so that $v\,w\to_\beta v'\,w'$ where $v\to_\beta v'$ and $w=w'$ or $v=v'$ and $w\to_\beta w'$. To handle this situation, we can use decomposition by means of rule $(\Rightarrow\mathrm E)$ then the induction hypothesis and assemble the necessary derivation again by rule $(\Rightarrow\mathrm E)$.
In the second case $v=\lambda x.v'$ and reduction follows the $\beta$-redex $(\lambda x.v')\,w\to_\beta v'[x:=w]$. Observe that the derivation of $\pmjudgment\Gamma{(\lambda x.v')\,w}\Bv$ must end with rule $(\Rightarrow\mathrm E)$
\[\frac{\pmjudgment\Gamma{\lambda x.v'}{R(\Av)\Rightarrow\Bv}\qquad\pmjudgment\Delta w\Av}{\pmjudgment{\Gamma\cup R(\Delta)}{(\lambda x.v')\,w}\Bv}\ (\Rightarrow\mathrm E)\]
for some relation $R$ and vector $\Av$. Then rule $(\Rightarrow\mathrm I)$ must be used to derive the judgment $\pmjudgment\Gamma{\lambda x.v'}{R(\Av)\Rightarrow\Bv}$, which gives derivation for $\pmjudgment{\Gamma,x:R(\Av)}{v'}\Bv$. We can now use the substitution lemma (cf. Lemma 23) to obtain $\pmjudgment{\Gamma\cup R(\Delta)}{v'[x:=w]}\Bv$, which is the desired conclusion.
\end{proof}
% SOURCE published pages17:9--17:10. Strict-system prerequisite and exact author projection proof.
\section*{Relation to the strict intersection type system}
As a technical tool to establish a relation with the traditional intersection type systems, we use the so-called strict system of van Bakel [22, Section 5]. This clarifies the relationship with the system of Barendregt, Coppo, Dezani-Ciancaglini [2], which is expressed in Theorem 30. In addition, this helps to establish metatheoretical properties, such as normalization.
\orig{Definition}{25 (Strict Intersection Type System [22, Definition 5.1])}
The rules of the strict system are slightly adapted to notation used in this paper.
\[\frac{A\in\sigma}{\pmstrictjudgment{\Sigma,x:\sigma}{x}{A}}\ (\cap\mathrm E)\qquad
\frac{\pmstrictjudgment\Sigma t{A_1},\ldots,\pmstrictjudgment\Sigma t{A_n}}{\pmstrictjudgment\Sigma t{\{A_1,\ldots,A_n\}}}\ (\cap\mathrm I)\]
\[\frac{\pmstrictjudgment{\Sigma,x:\sigma}t A}{\pmstrictjudgment\Sigma{\lambda x.t}{\sigma\to A}}\ (\to\mathrm I)\qquad
\frac{\pmstrictjudgment\Sigma t{\sigma\to B}\qquad\pmstrictjudgment\Sigma u\sigma}{\pmstrictjudgment\Sigma{t\,u}B}\ (\to\mathrm E)\]
\orig{Lemma}{27} If $\Gamma\vdash t:\Av$ then for each $i\in\{1,\ldots,\dimn(\Av)\}$ we have $\pmstrictjudgment{\Gamma_{[i]}}t{\Av_{[i]}}$.
\begin{proof}
Induction on $t$. Let $n:=\dimn(\Av)$ and assume $\Gamma\vdash t:\Av$.\par
\emph{Case $t=x$.} If $\Gamma\vdash x:\Av$ then it can be derived only by rule (Ax). Consequently, there is $(x:\sv)\in\Gamma$ such that $\Av\vin\sv$. Consequently, we can use rule $(\cap\mathrm E)$ to obtain $\pmstrictjudgment{\Gamma_{[i]}}x{\Av_{[i]}}$ for $i\in\{1,\ldots,n\}$.\par
\emph{Case $t=\lambda x.u$.} The judgment $\Gamma\vdash\lambda x.u:\Av$ can be derived only by rule $(\Rightarrow\mathrm I)$ and thus $\Av=\sv\Rightarrow\Bv$ for some $\sv$ and $\Bv$ and also $\Gamma,x:\sv\vdash u:\Bv$ is derivable. By the induction hypothesis, we obtain $\pmstrictjudgment{\Gamma_{[i]},x:\sv_{[i]}}u{\Bv_{[i]}}$ for each $i\in\{1,\ldots,n\}$. Then rule $(\to\mathrm I)$ gives us derivability of $\pmstrictjudgment{\Gamma_{[i]}}{\lambda x.u}{\sv_{[i]}\to\Bv_{[i]}}$. As $\sv_{[i]}\to\Bv_{[i]}=\Av_{[i]}$, we obtain the claim.\par
\emph{Case $t=u\,v$.} The judgment $\Gamma\vdash u\,v:\Av$ can be derived only by rule $(\Rightarrow\mathrm E)$. Consequently, $\Gamma=\Gamma'\cup R(\Delta)$ for some $\Gamma',\Delta$, and $R$ such that $R(\Bv)_{[i]}=\sv_{[i]}$ for $i\in\{1,\ldots,n\}$, and $\Gamma'\vdash u:\sv\Rightarrow\Av$ and $\Delta\vdash v:\Bv$ are derivable. By the induction hypothesis, we can derive
\begin{equation*}\begin{aligned}
\pmstrictjudgment{\Gamma'_{[i]}}u{\sv_{[i]}\to\Av_{[i]}}&\quad\text{for }i\in\{1,\ldots,n\},\\
\pmstrictjudgment{\Delta_{[j]}}v{\Bv_{[j]}}&\quad\text{for }j\in\{1,\ldots,\dimn(\Bv)\}.
\end{aligned}\tag{2}\end{equation*}
We have $\Gamma'(x)\vcup R(\Delta)(x)\vsub\Gamma(x)$ for $x\in\dom(\Gamma)$, and therefore $(\Delta(x))_{[j]}\subseteq(\Gamma(x))_{[i]}$ for $(j,i)\in R$. We can now apply the weakening lemma to derivations in (2) and obtain $\pmstrictjudgment{\Gamma_{[i]}}v{\Bv_{[j]}}$ for $(j,i)\in R$, as well as $\pmstrictjudgment{\Gamma_{[i]}}u{\sv_{[i]}\to\Av_{[i]}}$ for $i\in\{1,\ldots,n\}$. Using rule $(\cap\mathrm I)$ we derive $\pmstrictjudgment{\Gamma_{[i]}}v{\{\Bv_{[j]}\mid j\,R\,i\}}$. Application of rule $(\to\mathrm E)$ gives $\pmstrictjudgment{\Gamma_{[i]}}{u\,v}{\Av_{[i]}}$ for $i\in\{1,\ldots,n\}$.
\end{proof}
An important consequence of the lemma above is that many properties of van Bakel's system are inherited. In particular, system R has the strong normalization property for derivations without the empty vector, as van Bakel's system has [21, Theorem 7.5] for derivations without the empty type.
% SOURCE published pages17:12--17:13. Full author subformula and decision proof.
\section*{Decidability of Inhabitation in Bounded Dimension}
In this section, we give the main result of the present work: inhabitation in bounded dimension is decidable in 2-EXPTIME. This result is supplemented by the observation that for each fixed dimension inhabitation is decidable in EXPTIME.
\orig{Definition}{31 (Inhabitation in Bounded Dimension)} Given $\Gamma,\Av$, and $k$ is the judgment $\pmjudgment\Gamma t\Av$ derivable for some term $t$?
The proof of decidability is constructed in three steps. First, we establish inversion properties (Lemma 32 and Lemma 33) for abstraction and iterated application. Second, we formulate a subformula property (Lemma 36) for $\beta$-normal form derivations to restrict the inhabitant search space. Third, we give an alternating exponential space inhabitant search algorithm, obtaining a 2-EXPTIME upper bound on inhabitation (Theorem 37).
\orig{Lemma}{32} If $\pmjudgment\Gamma{\lambda x.t}\Bv$, then $\Bv=\sv\Rightarrow\Av$ for some $\sv$ and $\Av$, and $\pmjudgment{\Gamma,x:\sv}t\Av$.
\begin{proof}Case analysis regarding the last rule application.\end{proof}
\orig{Lemma}{33} If $\pmjudgment\Gamma{x\,u_1\ldots u_n}\Bv$ then there exists $\Av_1,\ldots,\Av_n,R_1,\ldots,R_n,\Delta_1,\ldots,\Delta_n$ such that
\begin{enumerate}
\item $R_i\subseteq\{1,\ldots,\dimn(\Av_i)\}\times\{1,\ldots,\dimn(\Bv)\}$ is left-total and $\dimn(\Av_i)\leq k$ for $i\in\{1,\ldots,n\}$.
\item $R_1(\Av_1)\Rightarrow\cdots\Rightarrow R_n(\Av_n)\Rightarrow\Bv\vin\Gamma(x)$.
\item $R_i(\Delta_i)\subseteq\Gamma$ for $i\in\{1,\ldots,n\}$.
\item $\pmjudgment{\Delta_i}{u_i}{\Av_i}$ for $i\in\{1,\ldots,n\}$.
\item $\pmjudgment\Gamma{x\,u_1\ldots u_i}{R_{i+1}(\Av_{i+1})\Rightarrow\cdots\Rightarrow R_n(\Av_n)\Rightarrow\Bv}$ for $i\in\{1,\ldots,n\}$.
\end{enumerate}
\begin{proof}
Induction on $n$ and case analysis regarding the last rule application. Left-totality of relations $R_1,\ldots,R_n$ for Property (1) is established using Lemma 17. Property (5) is established using Lemma 16.
\end{proof}
With each type or a set of types we associate a set of subformulae (types or sets of types), by the following Definition 34.
\orig{Definition}{34 (Subformula)}
\begin{align*}
\form(a)&:=\{a\},\\
\form(\sigma\to A)&:=\{\sigma\to A\}\cup\form(\sigma)\cup\form(A),\\
\form(\sigma)&:=\{\sigma\}\cup\bigcup_{A\in\sigma}\form(A).
\end{align*}
The notion of occurring subformulae is extended to vectors:
\[\form(\Av):=\bigcup_{i=1}^{\dimn(\Av)}\form(\Av_{[i]})\quad\text{and}\quad\form(\sv):=\bigcup_{i=1}^{\dimn(\sv)}\form(\sv_{[i]}).\]
The number of subformulae of a given type or set of types is polynomial in the number of nodes in its syntax tree, illustrated by the following example.
\orig{Example}{35} We have $\form(\{a,b\}\to c)=\{\{a,b\}\to c,\{a,b\},a,b,c\}$. In particular, the set $\{a,b\}$ and the types $a,b$ are subformulae of $\{a,b\}\to c$, but the sets $\{a\}$ and $\{b\}$ are not.
\orig{Lemma}{36 (Subformula Property)}
Assume that $\pmjudgment\Gamma t\Av$ such that $t$ is in $\beta$-normal form. Let $F:=\form(\Av)\cup\bigcup\{\form(\tv)\mid(y:\tv)\in\Gamma\}$ be the set of all subformulae of types and sets of types occurring in $\Gamma$ and $\Av$. There is a derivation of $\pmjudgment\Gamma t\Av$ such that for each judgment $\pmjudgment\Delta u\Bv$ in this derivation we have
\begin{enumerate}
\item for $(x:\sv)\in\Delta$ and $i\in\{1,\ldots,\dimn(\Delta)\}$ we have $\sv_{[i]}\in F\cup\{\emptyset\}$;
\item $\Bv_{[i]}\in F$ for $i\in\{1,\ldots,\dimn(\Bv)\}$.
\end{enumerate}
\begin{proof}
Induction on the size of $t$. There are two cases.\par
\emph{Case $t=\lambda x.u$:} Follows immediately from Lemma 32 and the induction hypothesis.\par
\emph{Case $t=x\,u_1\ldots u_n$:} There exist $\Av_1,\ldots,\Av_n,R_1,\ldots,R_n,\Delta_1,\ldots,\Delta_n$ satisfying properties of Lemma 33(1)--Lemma 33(5). Judgments $\pmjudgment\Gamma{x\,u_1\ldots u_i}{R_{i+1}(\Av_{i+1})\Rightarrow\cdots\Rightarrow R_n(\Av_n)\Rightarrow\Bv}$ (cf. Lemma 33(5)) satisfy properties (1) and (2). While subderivations $\pmjudgment{\Delta_i}{u_i}{\Av_i}$ for $i\in\{1,\ldots,n\}$ satisfy Property (2), they do not necessarily satisfy Property (1). By Lemma 33(3) type assumptions in $\Delta_i$ for $i\in\{1,\ldots,n\}$ may be strict subsets of subformulae in $F$. By Lemma 16 and Lemma 8(7) we can weaken the type assumptions by saturating the sets of types to subformulae in $F$. This results in derivation environments $\Delta'_1,\ldots,\Delta'_n$ and subderivations of judgments $\pmjudgment{\Delta'_i}{u_i}{\Av_i}$ for $i\in\{1,\ldots,n\}$ satisfying properties (1) and (2). We obtain the claim by the induction hypothesis and repeated application of rule $(\Rightarrow\mathrm E)$.
\end{proof}
% COUNTED STATEMENT 37

\orig{Theorem}{37} Given $\Gamma,\Av$, and $k$ it is decidable in 2-EXPTIME whether the judgment $\pmjudgment\Gamma t\Av$ is derivable for some term $t$.

% COUNTED PROOF 37
\begin{proof}
In order to decide whether $\pmjudgment\Gamma t\Av$ is derivable for some $t$, by subject reduction (Lemma 24) and normalization properties [21, Theorem 6.5 and Theorem 7.5] (combined with Lemma 27), it suffices to consider $t$ in $\beta$-normal form. In the following, we develop an AEXPSPACE decision procedure.
Let $F:=\form(\Av)\cup\bigcup\{\form(\tv)\mid(y:\tv)\in\Gamma\}$ be the set of all subformulae of types and sets of types occurring in $\Gamma$ and $\Av$, and let $N:=|F|$ denote its cardinality. By Lemma 36, we can consider a derivation of $\pmjudgment\Gamma t\Av$ such that for each judgment $\pmjudgment\Delta u\Bv$ in this derivation each coordinate of occurring type vectors can take at most $N+1$ distinct values (including the empty set for sets of types in environments). Therefore, the number of distinct $k$-dimensional vectors to consider in derivations is bounded by $(N+1)^k$.
The decision procedure works as follows. First, we non-deterministically guess whether $t$ is of shape $\lambda x.u$ for some $x$ and $u$, or of shape $x\,u_1\ldots u_n$ for some $x\in\dom(\Gamma)$ and some $u_1,\ldots,u_n$ where $n$ is bounded by the maximal nesting of the arrow type constructor to the right in types occurring in $\Gamma(x)$.
If $t$ is of shape $\lambda x.u$, then by Lemma 32 we have $\Av=\sv\Rightarrow\Av'$ for some $\sv$ and $\Av'$, and we recursively decide whether $\pmjudgment{\Gamma,x:\sv}u{\Av'}$ holds for some $u$.
If $t$ is of shape $x\,u_1\ldots u_n$, then we use Lemma 33, universally branching in $i\in\{1,\ldots,n\}$. In each branch $i$, we non-deterministically guess the necessary relation $R_i$ (bounded in cardinality by $k^2$), type vector $\Av_i$, environment $\Delta_i$, and recursively decide whether the judgment $\pmjudgment{\Delta_i}{u_i}{\Av_i}$ is derivable for some $u_i$.
By the pigeonhole principle, if a type environment contains more than $(N+1)^k$ variables, some must have identical assigned vectors. In the non-deterministic guess of the shape $\lambda x.u$, if $x$ is assigned the vector $\Gamma(y)$ for some $y\in\dom(\Gamma)$, we can equivalently guess $\lambda x.u[x:=y]$ instead and remove the assumption on $x$.
Overall, the space required by the alternating decision procedure is polynomial in $(N+1)^k$, yielding an AEXPSPACE, or equivalently, a 2-EXPTIME decision procedure.
\end{proof}
If we fix a dimension bound $k$, we obtain an EXPTIME inhabitation decision procedure.
% COUNTED STATEMENT 38 SAME RESULT FIXED PARAMETER

\orig{Theorem}{38} For any $k\in\mathbb N$, given $\Gamma$ and $\Av$ it is decidable in EXPTIME whether the judgment $\pmjudgment\Gamma t\Av$ is derivable for some term $t$.

% COUNTED PROOF 38
\begin{proof}
Similarly to the proof of Theorem 37, the number of distinct vectors in recursive invocation of the decision procedure is $(N+1)^k$, which is polynomial for any fixed $k$. This yields an APSPACE, or equivalently, an EXPTIME decision procedure.
\end{proof}

\begin{thebibliography}{99}
\bibitem[1]{ref1} H. P. Barendregt, W. Dekkers, and R. Statman. Lambda Calculus with Types. Perspectives in Logic, Cambridge University Press, 2013.
\bibitem[2]{ref2} Henk Barendregt, Mario Coppo, and Mariangiola Dezani-Ciancaglini. A filter lambda model and the completeness of type assignment. J. Symb. Log., 48(4):931–940, 1983. doi:10.2307/ 2273659.
\bibitem[3]{ref3} Antonio Bucciarelli, Delia Kesner, and Simona Ronchi Della Rocca. The inhabitation problem for non-idempotent intersection types. In TCS 2014, Proceedings of 8th IFIP TC 1/WG 2.2 International Conference on Theoretical Computer Science, Rome, Italy, September 1-3, 2014. Springer, 2014. doi:10.1007/978-3-662-44602-7\_26.
\bibitem[4]{ref4} Antonio Bucciarelli, Delia Kesner, and Daniel Ventura. Non-idempotent intersection types for the lambda-calculus. Log. J. IGPL, 25(4):431–464, 2017. doi:10.1093/JIGPAL/JZX018.
\bibitem[5]{ref5} M. Coppo, M. Dezani-Ciancaglini, and B. Venneri. Functional Characters of Solvable Terms. Zeitschrift für mathematische Logik und Grundlagen der Mathematik, pages 45–58, 1981.
\bibitem[6]{ref6} Daniel de Carvalho. Execution time of \ensuremath{\lambda}-terms via denotational semantics and intersection types. Math. Struct. Comput. Sci., 28(7):1169–1203, 2018. doi:10.1017/S0960129516000396.
\bibitem[7]{ref7} Andrej Dudenhefner and Jakob Rehof. Intersection type calculi of bounded dimension. POPL 2017, SIGPLAN Not., 52(1):653–665, January 2017. doi:10.1145/3093333.3009862.
\bibitem[8]{ref8} Andrej Dudenhefner and Jakob Rehof. Typability in bounded dimension. In 32nd Annual ACM/IEEE Symposium on Logic in Computer Science (LICS), pages 1–12, 2017. doi: 10.1109/LICS.2017.8005127.
\bibitem[9]{ref9} Andrej Dudenhefner, Christoph Stahl, Constantin Chaumet, Felix Laarmann, and Jakob Rehof. Finite Combinatory Logic with Predicates. In Delia Kesner, Eduardo Hermo Reyes, and Benno van den Berg, editors, 29th International Conference on Types for Proofs and Programs (TYPES 2023), volume 303 of Leibniz International Proceedings in Informatics (LIPIcs), pages 2:1–2:22, Dagstuhl, Germany, 2024. Schloss Dagstuhl – Leibniz-Zentrum für Informatik. doi:10.4230/LIPIcs.TYPES.2023.2.
\bibitem[10]{ref10} Andrej Dudenhefner and Paweł Urzyczyn. Kripke semantics for intersection formulas. ACM Trans. Comput. Log., 22(3):15:1–15:16, 2021. doi:10.1145/3453481.
\bibitem[11]{ref11} Jonathan Frankle, Peter-Michael Osera, David Walker, and Steve Zdancewic. Example-directed synthesis: a type-theoretic interpretation. POPL 2026, ACM Sigplan Notices, 51(1):802–815, 2016. doi:10.1145/2837614.2837629.
\bibitem[12]{ref12} J. Roger Hindley. The simple semantics for coppo-dezani-sallé types. In Mariangiola Dezani- Ciancaglini and Ugo Montanari, editors, International Symposium on Programming, 5th Colloquium, Torino, Italy, April 6-8, 1982, Proceedings, volume 137 of Lecture Notes in Computer Science, pages 212–226. Springer, 1982. doi:10.1007/3-540-11494-7\_15.
\bibitem[13]{ref13} D. Leivant. Polymorphic Type Inference. In POPL 1983, Proceedings of the 10th ACM Symposium on Principles of Programming Languages, pages 88–98, 1983. doi:10.1145/ 567067.567077.
\bibitem[14]{ref14} Luigi Liquori and Simona Ronchi Della Rocca. Intersection-types à la church. Inf. Comput., 205(9):1371–1386, 2007. doi:10.1016/J.IC.2007.03.005.
\bibitem[15]{ref15} Elaine Pimentel, Simona Ronchi Della Rocca, and Luca Roversi. Intersection types from a proof-theoretic perspective. Fundam. Informaticae, 121(1-4):253–274, 2012. doi:10.3233/ FI-2012-778.
\bibitem[16]{ref16} Nadia Polikarpova, Ivan Kuraj, and Armando Solar-Lezama. Program synthesis from polymorphic refinement types. PLDI 2016, ACM SIGPLAN Notices, 51(6):522–538, 2016. doi:10.1145/2908080.2908093.
\bibitem[17]{ref17} Simona Ronchi Della Rocca and Luca Roversi. Intersection logic. In Laurent Fribourg, editor, Computer Science Logic, 15th International Workshop, CSL 2001., volume 2142 of Lecture Notes in Computer Science, pages 414–428. Springer, 2001. doi:10.1007/3-540-44802-0\_29.
\bibitem[18]{ref18} Simona Ronchi Della Rocca, Alexis Saurin, Yiorgos Stavrinos, and Anastasia Veneti. Intersec- tion logic in sequent calculus style. In Elaine Pimentel, Betti Venneri, and Joe B. Wells, editors, Proceedings Fifth Workshop on Intersection Types and Related Systems, ITRS 2010, Edinburgh, U.K., 9th July 2010., volume 45 of EPTCS, pages 16–30, 2010. doi:10.4204/EPTCS.45.2.
\bibitem[19]{ref19} P. Urzyczyn. Inhabitation of Low-Rank Intersection Types. In TLCA 2009, Proceedings of Typed Lambda Calculus and Applications, volume 5608 of LNCS, pages 356–370. Springer, 2009. doi:10.1007/978-3-642-02273-9\_26.
\bibitem[20]{ref20} Steffen van Bakel. Complete restrictions of the intersection type discipline. Theoretical Computer Science, 102(1):135–163, 1992. doi:10.1016/0304-3975(92)90297-S.
\bibitem[21]{ref21} Steffen van Bakel. Cut-elimination in the strict intersection type assignment system is strongly normalizing. Notre Dame J. Formal Log., 45(1):35–63, 2004. doi:10.1305/NDJFL/1094155278.
\bibitem[22]{ref22} Steffen van Bakel. Strict intersection types for the lambda calculus. ACM Comput. Surv., 43(3):20:1–20:49, 2011. doi:10.1145/1922649.1922657.
\bibitem[23]{ref23} J. B. Wells, Allyn Dimock, Robert Muller, and Franklyn A. Turbak. A calculus with polymorphic and polyvariant flow types. J. Funct. Program., 12(3):183–227, 2002. doi: 10.1017/S0956796801004245.
\bibitem[24]{ref24} Joe B. Wells and Christian Haack. Branching types. In Daniel Le Métayer, editor, Programming Languages and Systems, 11th European Symposium on Programming, ESOP 2002, held as Part of the Joint European Conference on Theory and Practice of Software, ETAPS 2002, Grenoble, France, April 8-12, 2002, Proceedings, volume 2305 of Lecture Notes in Computer Science, pages 115–132. Springer, 2002. doi:10.1007/3-540-45927-8\_9.
\end{thebibliography}
\end{document}
